\chapter{Mô tả hệ thống mục tiêu}
\label{ch:hethong}

\section{Chức năng và luồng xử lý}
Đơn hàng được tạo (\code{CREATED}), chuyển sang chờ cổng thanh toán (\code{PENDING}). Cổng có
thể xác nhận thu tiền (\code{SUCCEEDED}), báo thất bại (\code{FAILED}) hoặc không phản hồi
(\code{UNKNOWN} — timeout); đơn timeout phải được đối soát (\code{RECONCILING}) để biết kết
quả thật. Chỉ đơn đã thu tiền mới được giao (\code{SHIPPED}) hoặc hoàn tiền
(\code{REFUNDING} $\rightarrow$ \code{REFUNDED}). Song song, worker xử lý webhook và worker
đối soát cùng ghi vào sổ cái, được bảo vệ bởi một khóa. Kết quả từ cổng về qua
\emph{webhook} — điểm nhận dữ liệu từ bên ngoài quan trọng nhất (Hình~\ref{fig:kientruc}).

\begin{figure}[H]
  \centering
  \resizebox{\linewidth}{!}{% Hình: kiến trúc luồng thanh toán và trust boundary
\begin{tikzpicture}[node distance=7mm and 9mm]
  \node[box] (kh) {Khách hàng\\(trình duyệt)};
  \node[box,right=of kh] (shop) {Order Service\\(tạo đơn)};
  \node[box,right=of shop] (pay) {Payment Service\\\code{parse\_webhook\_signature}\\\code{calculate\_total\_payment}};
  \node[box,below=12mm of pay] (ful) {Fulfilment\\(giao hàng)};
  \node[box,right=16mm of pay,fill=black!5] (gw) {Cổng thanh toán\\(bên thứ ba)};

  \draw[arr] (kh) -- node[above,font=\scriptsize]{đặt hàng} (shop);
  \draw[arr] (shop) -- node[above,font=\scriptsize]{start} (pay);
  \draw[arr] (pay.north east) to[bend left=20] node[above,font=\scriptsize]{yêu cầu thu} (gw.north west);
  \draw[arr,bad,thick] (gw.south west) to[bend left=20] node[below,font=\scriptsize,text=bad]{webhook} (pay.south east);
  \draw[arr] (pay) -- node[right,font=\scriptsize]{ship (chỉ khi paid)} (ful);

  \begin{scope}[on background layer]
    \node[draw=bad,dashed,thick,rounded corners,inner sep=4mm,fit=(shop)(pay)(ful),
          label={[font=\scriptsize,text=bad]above:Trust boundary — hệ thống nội bộ}] {};
  \end{scope}
\end{tikzpicture}
}
  \caption{Kiến trúc luồng thanh toán; đường đỏ là dữ liệu vượt trust boundary.}
  \label{fig:kientruc}
\end{figure}

\section{Bề mặt tấn công (attack surface)}
\begin{table}[H]
\centering\small
\caption{Các điểm nhận dữ liệu từ bên ngoài và rủi ro tương ứng.}
\label{tab:attacksurface}
\begin{tabularx}{\linewidth}{@{}l >{\raggedright\arraybackslash}X >{\raggedright\arraybackslash}X l@{}}
\toprule
Điểm vào & Dữ liệu & Hàm xử lý & Rủi ro \\
\midrule
Webhook & Chữ ký HMAC-SHA256 (64 ký tự hex) & \code{parse\_webhook\_signature} & CWE-120 \\
Webhook & Header token & \code{webhook\_bypass\_token} & CWE-798 \\
API hỗ trợ viên & Mã PIN ghi đè & \code{support\_override\_pin} & CWE-798 \\
Tạo đơn & Mã đơn hàng (20 ký tự) & \code{parse\_order\_id} & CWE-120 \\
Webhook & Số tiền, phí vận chuyển & \code{calculate\_total\_payment} & CWE-190 \\
API hoàn tiền & Số tiền hoàn, tổng tiền & \code{refund\_ratio} & CWE-369 \\
Nội bộ & Bản ghi audit, phiên, token thẻ & \code{write\_audit\_log}, \code{close\_payment\_session}, \code{release\_card\_token\_buffer} & CWE-401/415/416 \\
\bottomrule
\end{tabularx}
\end{table}

\section{Mô hình mối đe dọa (STRIDE)}
\begin{table}[H]
\centering\small
\caption{Phân tích STRIDE cho luồng thanh toán.}
\label{tab:stride}
\begin{tabularx}{\linewidth}{@{}l X l l@{}}
\toprule
Mối đe dọa & Kịch bản & Thuộc tính bị vi phạm & Yêu cầu \\
\midrule
Spoofing & Gửi webhook giả ``thanh toán thành công'' bằng token bí mật & Authenticity & REQ-11 \\
Tampering & Sửa số tiền trong payload làm tổng bị tràn số & Integrity & REQ-9 \\
Repudiation & Mất bản ghi audit khi worker crash & Non-repudiation & NFR-3 \\
Info.\ Disclosure & Tràn bộ đệm chữ ký làm lộ dữ liệu stack & Confidentiality & REQ-5 \\
Denial of Service & Đơn treo mãi ở trạng thái timeout; hoàn tiền không kết thúc & Availability & REQ-2, REQ-13 \\
Elevation of Priv. & PIN ghi đè hardcode của hỗ trợ viên & Authorization & REQ-11 \\
\bottomrule
\end{tabularx}
\end{table}

\section{Đặc tả yêu cầu}
Đặc tả viết theo template Excel 10 sheet của học phần và được sinh tái lập bằng
\code{scripts/build\_payment\_specs.py} thành hai file
\code{payment\_gateway\_v0\_spec.xlsx} và \code{payment\_gateway\_v1\_spec.xlsx} trong thư mục
\code{dataset/specs\_excel/}.
Hai phiên bản dùng chung yêu cầu, thuộc tính và test case; chỉ khác sheet
\code{MoHinhTrangThai} và cột \code{Loi} của \code{HamSinhMa}.

\begin{table}[H]
\centering\small
\caption{Yêu cầu trong sheet \code{YeuCauChucNang} (rút gọn).}
\label{tab:req}
\begin{tabularx}{\linewidth}{@{}l l l X@{}}
\toprule
ID & Danh mục & Ưu tiên & Nội dung \\
\midrule
REQ-1 & safety & HIGH & Chỉ giao hàng khi đơn đã được xác nhận thu tiền \\
REQ-2 & liveness & HIGH & Đơn timeout cuối cùng phải được đối soát xong \\
REQ-3 & safety & HIGH & Hai worker không đồng thời giữ khóa sổ cái \\
REQ-4 & liveness & HIGH & Không deadlock khi tranh chấp khóa hoặc khi hoàn tiền \\
REQ-5 & memory & HIGH & \code{parse\_webhook\_signature} không tràn bộ đệm \\
REQ-6 & memory & MEDIUM & \code{write\_audit\_log} giải phóng bộ nhớ mỗi bản ghi \\
REQ-7 & memory & HIGH & \code{close\_payment\_session} không free hai lần \\
REQ-8 & memory & HIGH & \code{release\_card\_token\_buffer} không dùng sau khi free \\
REQ-9 & safety & HIGH & Tổng tiền hàng + phí không tràn int32 \\
REQ-10 & functional & MEDIUM & \code{refund\_ratio} không lỗi khi tổng tiền bằng 0 \\
REQ-11 & safety & HIGH & Không có mã bí mật hardcode vượt qua xác thực \\
REQ-12 & safety & HIGH & Không ghi nhận thu tiền hai lần khi webhook bị gửi lại \\
REQ-13 & liveness & MEDIUM & Mọi yêu cầu hoàn tiền cuối cùng phải hoàn tất \\
REQ-14 & memory & MEDIUM & \code{parse\_order\_id} chỉ nhận mã đơn đúng độ dài \\
\bottomrule
\end{tabularx}
\end{table}

\begin{table}[H]
\centering\small
\caption{Quy mô đặc tả (v0 / v1).}
\label{tab:quymo}
\begin{tabular}{@{}lr@{}}
\toprule
Thành phần & Số lượng \\
\midrule
Yêu cầu chức năng / phi chức năng / ràng buộc & 14 / 6 / 6 \\
Trạng thái / chuyển tiếp của mô hình Kripke & 14 / 23 \quad (v1: 12 / 18) \\
Công thức CTL / cặp loại trừ lẫn nhau & 8 / 5 \\
Hàm sinh mã (C / Python) & 6 / 3 \\
Kiểm tra số học (Z3) & 6 \\
Test case truy vết & 12 \\
Số dòng mã C sinh ra (v0 + v1) & 338 \\
\bottomrule
\end{tabular}
\end{table}

Đặc tả đáp ứng checklist của Hướng dẫn tạo spec Excel: đủ 10 sheet, mỗi sheet dạng bảng
$\geq 5$ dòng, phủ 4 danh mục yêu cầu, có cụm khóa \code{free/holds/conflict} với nhãn
\code{violation}, công thức sống còn dạng $\mathbf{AG}(p \rightarrow \mathbf{AF}\,q)$, đủ 4 loại
lỗi bộ nhớ C cùng hai loại hàm Python \code{divide} và \code{magic\_value}; kích thước bộ đệm
khớp các ràng buộc CON-1 đến CON-3.
